% Ninja -- Round 2 demo video script
% Compile with: pdflatex demo_script.tex
\documentclass[11pt,a4paper]{article}

\usepackage[T1]{fontenc}
\usepackage[utf8]{inputenc}
\usepackage[margin=2.1cm]{geometry}
\usepackage{lmodern}
\usepackage{xcolor}
\usepackage{enumitem}
\usepackage{listings}
\usepackage{booktabs}
\usepackage{titlesec}
\usepackage[hidelinks]{hyperref}

% ---------------------------------------------------------------- palette
\definecolor{ink}{HTML}{14181F}
\definecolor{dim}{HTML}{5B6B7A}
\definecolor{accent}{HTML}{2F6FB5}
\definecolor{critical}{HTML}{C0392B}
\definecolor{healthy}{HTML}{2F7D46}
\definecolor{rulegrey}{HTML}{D3D9E0}
\definecolor{shell}{HTML}{F4F6F8}

\color{ink}

% ---------------------------------------------------------------- code
\lstdefinestyle{shellstyle}{
  basicstyle=\ttfamily\small,
  backgroundcolor=\color{shell},
  frame=leftline,
  framerule=1.6pt,
  rulecolor=\color{accent},
  framesep=7pt,
  xleftmargin=10pt,
  breaklines=true,
  columns=fullflexible,
  keepspaces=true,
  aboveskip=8pt,
  belowskip=8pt,
}
\lstset{style=shellstyle}

% ---------------------------------------------------------------- beats
\newcommand{\beat}[3]{%
  \vspace{9pt}%
  \noindent{\color{rulegrey}\rule{\linewidth}{0.6pt}}\\[3pt]
  \noindent{\ttfamily\bfseries\color{accent}#1}\quad{\bfseries\large #2}%
  \if\relax\detokenize{#3}\relax\else\\[1pt]{\small\color{dim}#3}\fi
  \par\vspace{5pt}%
}
\newcommand{\visual}[1]{\noindent{\small\textbf{\color{dim}VISUAL}\quad #1}\par\vspace{3pt}}
\newcommand{\action}[1]{\noindent{\small\textbf{\color{dim}ON SCREEN}\quad #1}\par\vspace{3pt}}
\newcommand{\vo}[1]{\vspace{2pt}\noindent\begin{minipage}{\linewidth}%
  \small\textbf{\color{dim}VO}\quad\itshape #1\end{minipage}\par\vspace{4pt}}
\newcommand{\warn}[1]{\vspace{3pt}\noindent\fcolorbox{critical}{white}{%
  \begin{minipage}{\dimexpr\linewidth-2\fboxsep-2\fboxrule}%
  \small\textbf{\color{critical}CAREFUL}\quad #1\end{minipage}}\par\vspace{4pt}}

\titleformat{\section}{\large\bfseries}{}{0pt}{}
\setlength{\parindent}{0pt}

% ================================================================
\begin{document}

\begin{center}
  {\LARGE\bfseries Ninja — Round 2 Demo Video Script}\\[5pt]
  {\color{dim}Target runtime $\approx$ 4:10 \quad|\quad screen recording + voiceover}\\[3pt]
  {\small\color{dim}Every figure below is taken from the built prototype:
  \texttt{reports/holdout\_report.md} and \texttt{reports/holdout\_results.json}.
  Nothing is invented for the video.}
\end{center}

\vspace{6pt}
\noindent\fcolorbox{rulegrey}{shell}{%
\begin{minipage}{\dimexpr\linewidth-2\fboxsep-2\fboxrule}
\small
\textbf{Three corrections from the previous draft.} Each was a factual error against the
actual build.
\begin{enumerate}[leftmargin=14pt,itemsep=2pt,topsep=3pt]
  \item \textbf{The uninstrumented station is 34, not 12.} Station 12 \emph{has}
  sensors. The 14 unseen stations are 3, 5, 7, 11, 14, 17, 22, 25, 27, 29, 34,
  36, 37, 39. Naming station 12 as sensorless on camera is the kind of error a
  judge who opened \texttt{config/line.yaml} would catch.
  \item \textbf{\texttt{run\_demo.py} does not print the holdout numbers.} It
  prints \textsc{demo scorecard (tuning split)} — development figures, and it
  says so on screen. The sealed numbers come from
  \texttt{python -m eval.run\_holdout}. The old script narrated holdout results
  over tuning output. Now split into two beats.
  \item \textbf{No supply-chain dataset is currently shipped.} \texttt{data/} has
  no CSV, so \texttt{config/line.yaml} records \texttt{source: defaults}. Claiming
  it \emph{is} calibrated against real data overstates it. Corrected wording is
  in the honesty beat.
\end{enumerate}
\end{minipage}}

% ================================================================
\beat{0:00–0:18}{Cold open — show it run before explaining it}{}
\visual{Terminal, empty. Type the command live.}
\begin{lstlisting}
pip install -r requirements.txt
python run_demo.py
\end{lstlisting}
\vo{``This is Ninja — a digital twin of a vehicle assembly line. This is a
cold start, on a clean checkout. No GPU, no network, no API keys.''}
\action{Let the scorecard print. Roughly 12 seconds. Do not cut away.}
\vo{``Four scenarios: simulator, counterfactual, twin, scoring. Twelve seconds.''}

% ================================================================
\beat{0:18–0:50}{The claim — and the sealed numbers}{This is the beat the old script got wrong.}
\warn{The \texttt{run\_demo.py} output is headed \textsc{demo scorecard (tuning split)}
and explicitly labels itself development figures. Do not read the holdout numbers
over it. Run the holdout scorer as a second command — it is five seconds of
screen time and it is the honest version.}
\begin{lstlisting}
python -m eval.run_holdout
\end{lstlisting}
\visual{The \textsc{holdout scorecard} table printing, per scenario.}
\vo{``That was the development set. This is the one that counts — nine scenarios
the system has never been tuned against, scored once, against thresholds that
were frozen before it ever saw them.''}
\vo{``Six faults out of six detected. A forming bottleneck called a median of
\textbf{37 minutes} before the line actually backs up. \textbf{99\%} of units
built through a drifted process flagged before they reach final inspection —
387 of 390. And \textbf{5.5} false alarms per seven-and-a-half-hour shift,
across 40 stations.''}
\vo{``Every one of those is reproducible from a clean checkout. The rest of this
video is why you should believe them rather than just take them.''}

% ================================================================
\beat{0:50–1:20}{The problem, in one breath}{}
\visual{The Supervisor tab, static. The 40-station line view.}
\vo{``Two things go wrong on a real line. A station drifts slower, and nobody
notices until a queue has already formed downstream. And a torque tool trends
out of tolerance, producing cars that pass every local check and fail dozens of
units later at final inspection.''}
\vo{``Dashboards report both after the fact. And there is a harder problem
underneath: on this line only 26 of 40 stations have sensors at all. The other
14 emit nothing per unit. Buffer levels are never reported anywhere.''}
\action{Point at the hollow diamonds in the line view.}
\vo{``Those hollow diamonds are the stations with no sensors. Notice their
uncertainty bands are visibly wider than the filled circles. That is the twin
telling you how much it does not know.''}

% ================================================================
\beat{1:20–2:05}{Why the numbers are trustworthy}{The counterfactual, then two live proofs.}
\visual{The architecture diagram from the README: simulator $\rightarrow$
events.jsonl $\rightarrow$ twin $\rightarrow$ eval, ground truth branching off separately.}
\vo{``Here is the part that actually matters. Every scenario runs twice from the
same random seed — once with a fault injected, once without — and the two runs
share their per-unit random draws. So any difference between them is caused by
the fault and nothing else.''}
\vo{``That is how we know exactly when a bottleneck would have formed. Not an
estimate. A computed fact. Every lead-time number is measured against it.''}
\vo{``And the twin never sees any of that. It reads one file: the event stream.
We do not just assert that — we test it.''}
\begin{lstlisting}
pytest tests/test_isolation.py -v
\end{lstlisting}
\vo{``Three tests. The twin runs in a separate process, in a folder where the
ground-truth file does not exist. And the third one asserts its output is
byte-identical whether or not ground truth is sitting on disk beside it. A grep
for imports would prove nothing; this does.''}
\begin{lstlisting}
python -m eval.freeze_thresholds --verify
\end{lstlisting}
\vo{``And the thresholds are sealed with a hash. If the config had changed after
freezing, the holdout scorer would refuse to run at all — not warn, refuse.''}

% ================================================================
\beat{2:05–3:00}{Live perturbation — the anti-scripted-demo proof}{Do this as one continuous take.}
\visual{Streamlit app, Validation tab, Live perturbation panel.}
\vo{``This is the part a scripted demo cannot fake. We are going to inject a
fault live, right now.''}
\action{The panel already defaults to \textbf{S34 — FINA-34 Exhaust Hang}. A blue
callout appears stating the station is uninstrumented. Leave the defaults: drift,
severity 30, onset 210 min.}
\vo{``Station 34. One of the 14 stations on this line with no sensors at all —
the app is telling you that itself, right there. The twin has no direct reading
from this station. Watch.''}
\action{Click \textbf{Run scenario}. It takes 1--3 seconds. Do not cut.}
\vo{``That just ran a fresh simulation, ran its same-seed counterfactual, fed
only the event stream to the twin, and scored the result with the same matching
rule as the holdout report.''}
\visual{The result row: \textsc{detected} yes \quad \textsc{lead time} +54 min
\quad \textsc{units lost} 21.}
\vo{``Detected — 54 minutes before the buffer actually saturated. On a station
that has no sensor on it.''}
\vo{``L1 caught it from the timing of its instrumented neighbours alone. L2
projected forward and named the specific downstream station that would starve.
And the confidence on that alert is \emph{medium}, not high, with a band roughly
three times wider than the same fault at an instrumented station.''}
\action{Change the target station or drag the severity slider. Re-run.}
\vo{``Change the fault and the answer changes with it. Nothing here is
precomputed.''}

% ================================================================
\beat{3:00–3:30}{The trade-off, shown not told}{}
\visual{The master sensitivity slider in the sidebar, and the threshold-sweep
curve on the Validation tab.}
\vo{``Every detector like this has one trade-off: catch things earlier, or raise
fewer false alarms. We do not just mention it — we plot it.''}
\action{Drag the sensitivity slider. The false-alarms-per-shift figure in the
sidebar updates live, and the marker moves along the sweep curve.}
\vo{``Turn sensitivity up and false alarms climb — at maximum it is over a
hundred per shift and precision collapses. Turn it down and you lose recall.
The 5.5 we report is the operating point we chose, and you can see exactly what
we traded to get it.''}
\vo{``That curve is real measured data from the sealed scenarios, not an
illustration.''}

% ================================================================
\beat{3:30–3:55}{The honesty beat — say the limitations first}{}
\visual{The ``Honest limitations'' section of the README.}
\vo{``Three things we want to say before anyone asks.''}
\vo{``First: all of this runs on simulated data, deliberately — because
simulation is the only way to know ground truth well enough to score yourself
honestly. The calibration script fits inbound-delay and defect-rate parameters
from a real supply-chain dataset when one is supplied. We are shipping without
one, so it is currently running on documented defaults: a 0.5\% base defect rate,
the order of magnitude of the public Bosch production-line dataset, and drift
shapes following AI4I 2020. The config records which of those it used.''}
\vo{``Second: the false-alarm rate comes from three control runs, so the
confidence interval is wide — 0.1 to 10.9. We report the interval rather than
hiding it.''}
\vo{``Third: our bottleneck layer is accurate rather than early. The early
warning comes from drift detection. Propagation's job is naming where and what,
not adding lead time on its own.''}

% ================================================================
\beat{3:55–4:10}{Close}{}
\visual{The README headline table, then the holdout report.}
\vo{``Tuned on one set of scenarios, seeds 1000 to 1999. Scored once, on another
set it never saw, seeds 9000 to 9999. The ranges are enforced at config load
time, so a holdout scenario cannot be an accidental replay of a tuned one.''}
\vo{``51 tests. Every number in this video is in that report, and every one is
reproducible from a clean checkout. That is Ninja.''}
\action{Cut to black / logo card.}

% ================================================================
\vspace{10pt}
\section*{Trim guide — if you need 3:00}
\begin{enumerate}[leftmargin=16pt,itemsep=3pt]
  \item Cut the trade-off beat (3:00--3:30). Replace with one line: ``we also
  tuned this against a visible false-alarm trade-off, plotted in the repo.''
  \textbf{Saves $\approx$30s.}
  \item Jump-cut the install in the cold open --- ``\ldots and a few seconds
  later'' --- rather than watching pip run. \textbf{Saves $\approx$8s.}
  \item Compress the honesty beat to one sentence: ``everything here is simulated
  by design, it currently runs on documented defaults rather than a supplied
  dataset, and our confidence intervals reflect a small control sample --- all
  three stated plainly in the README.'' \textbf{Saves $\approx$15s.}
  \item Drop the \texttt{freeze\_thresholds --verify} command, keep the sentence.
  \textbf{Saves $\approx$7s.}
\end{enumerate}

\vspace{4pt}
\noindent\textbf{Never cut:} the counterfactual explanation, the live perturbation
on the uninstrumented station, and the closing ``tuned on one set, scored on
another'' line. Those three are what separate this from every other submission.

% ================================================================
\section*{Recording notes}
\begin{itemize}[leftmargin=16pt,itemsep=3pt]
  \item \textbf{The perturbation segment must be one continuous take.} A cut
  there is the first thing a sceptical judge will question.
  \item Start the app before recording --- first load simulates and runs the twin,
  which takes a few seconds. Once warm it is cached and instant.
  \item Pre-run the isolation test off camera so you know it passes, then run it
  live anyway.
  \item If \texttt{run\_demo.py} takes its full $\approx$12 seconds, let it
  breathe. A visibly working cold start is worth more than saving five seconds.
  \item The console is dark-themed; record on a dark desktop background so the
  cut to the terminal is not jarring.
\end{itemize}

% ================================================================
\section*{Figures used, and where they come from}
\begin{center}
\small
\begin{tabular}{@{}lll@{}}
\toprule
\textbf{Figure} & \textbf{Value} & \textbf{Source} \\
\midrule
Median lead time (queue-forming, n=4) & +37.4 min & holdout\_results.json \\
Defect containment (n=2)              & 99.2\%, 387/390 & holdout\_results.json \\
Faults detected                       & 6 / 6 & holdout\_results.json \\
Alert precision                       & 74.2\% & holdout\_results.json \\
False alarms per shift (n=3 controls) & 5.50, CI 0.09--10.91 & holdout\_results.json \\
Uninstrumented holdout fault (S34)    & +34 min early & holdout\_report.md \\
Live perturbation, S34 @ 30\%         & +54 min, 21 units lost & app, live \\
Instrumented coverage                 & 26 of 40 (65\%) & config/line.yaml \\
Takt                                  & 76.0 s & config/line.yaml \\
Test count                            & 51 & pytest \\
Cold-start demo runtime               & $\approx$12 s & run\_demo.py \\
\bottomrule
\end{tabular}
\end{center}

\vspace{4pt}
{\small\color{dim}\noindent The two lead-time figures are different measurements
and both are real: \textbf{+34 min} is the sealed holdout result for the
uninstrumented station; \textbf{+54 min} is what the live panel produces at its
default settings, on a shorter 8-hour horizon. Use each only where this script
places it.}

\end{document}
